\chapter{Gossip Models}

Introduciamo una nuova classe di modelli distribuiti, i modelli GOSSIP. Dato un grafo $G = (V, E)$, con $|V| = n\ \&\ |E| = m$, denotiamo $\forall v \in V$, il vicinato di $v$ come $N(v)$ e $\delta(v) = |N(v)|$ la cardinalità del suo vicinato incluso se stesso. (Attenzione, diverso dal suo grado che è $\delta(v) - 1$).

\begin{definizione}[The $\mathcal{PUSH (PULL)}$ Model]
    Dato un grafo (non diretto) $G=(V, E)$, il modello di comunicazione uniforme $\mathcal{PUSH (PULL)}$ funziona nel mondo \textit{sincrono}, a \textit{round discreti} come segue. Ad ogni round $t = 0,1, \dots$, ciascun nodo $v \in V$ sceglie u.a.r uno dei suoi vicini $u \in_{u} N(v)$:
    \begin{description}
        \item[$\mathcal{PUSH}$] $v$ trasmette un messaggio $M$ a $u$, che lo riceverà entro il time slot $t$.
        \item[$\mathcal{PULL}$] $v$ si può far trasmettere da $u$ un qualsiasi messaggio $M$, che lo riceverà entro il time slot $t$.
    \end{description}
\end{definizione}

\section{Proprietà}

Si osserva che ad ogni istante $t$ tutte le operazioni di \textit{push} o \textit{pull} generano un \textbf{grafo diretto delle comunicazioni} $G_t = {V, E_t}$. Consideriamo il protocollo $\mathcal{PUSH}$, esiste un arco diretto $(u, v) \in E_t$ se $u$ fa un operazione di push su $v$. Analogamente per il protocollo $\mathcal{PULL}$, esiste l'arco diretto $(v, u) \in E_t$ se $u$ fa un operazione di pull su $v$. Osservare inoltre che il grafo $G_t$ risulta \textbf{sparso}, ovvero con $n$ archi, perché ci sarà un arco per ogni operazione di \textit{push} o \textit{pull} fatta dai nodi.

\begin{lemma}
    Valgono i seguenti fatti:
    \begin{enumerate}
        \item Nel modello $\mathcal{PUSH}$, ad ogni round, il numero atteso di operazioni di \textit{push} ricevute da ciascun nodo è 1 ed è $O(\log n)$ w.h.p.
        \item Nel modello $\mathcal{PULL}$, ad ogni round, il numero atteso di operazioni di \textit{pull} ricevute da ciascun nodo è 1 ed è $O(\log n)$ w.h.p.
    \end{enumerate}
    \begin{proof}
        Dimostriamo formalmente solo il primo fatto, per il secondo la dimostrazione è analoga. Consideriamo come grafo una \textit{clique} $K_n$. Poichè questo è il grafo più denso, i risultati varranno anche per grafi non completi.

        Possiamo modellare questo processo come \textit{balls into bins}, dove le $n$ palline e gli $n$ sono i $n$ nodi del grafo.  
        
        \begin{center}
            \begin{tikzpicture}[
                vertex/.style = {circle, draw, minimum size=8mm},
                arrow/.style = {->, >={Stealth[scale=1.2]}, thick, blue},
                % Stile per le etichette dei gruppi
                group label/.style = {font=\bfseries\small, align=center}
                ]

                % --- Creazione Nodi ---
                \foreach \i in {1,...,4} {
                    \node[vertex] (a\i) at (0, {-\i*1.5}) {};
                    \node[vertex] (b\i) at (4, {-\i*1.5}) {};
                }

                % --- Etichette Gruppi (V1 e V2) ---
                % Posizionate sopra il primo nodo di ogni colonna
                \node[group label, node distance=0.5cm, above=of a1] ($V_1$) {V1 \\ (nodi pusher)};
                \node[group label, node distance=0.5cm, above=of b1] ($V_2$) {V2 \\ (nodi receiver)};

                % --- Etichette u e v ---
                \node[left=3mm, blue, font=\boldmath] at (a2) {u};
                \node[right=3mm, blue, font=\boldmath] at (b2) {v};

                % --- Frecce ---
                \draw[arrow] (a2) -- (b2); % Da u a v
                \draw[arrow] (a1) -- (b2); % Altra entrante in v
                \draw[arrow] (a3) -- (b2); % Altra entrante in v

            \end{tikzpicture}
        \end{center}
        Dalla figura, osserviamo che il grado uscente di ciascun nodo è 1 ed è un valore \textit{deterministico}. Quello che ci interessa invece determinare è i grado entrante, che è una v.a. definita nel segente modo: Sia $v \in V_2 = V$ un generico nodo fisato e si consideri la variabile aleatoria binaria
        \[
        X^{(t)}_{u,v} = 
            \begin{cases*}
            1 \texttt{ if } u \xrightarrow{\texttt{push}} v \\
            0 \text{ altrimenti }
            \end{cases*}
        \]
        Quindi, $X^{(t)}_v = \sum_{u \in V_1} X^{(t)}_{u,v}$ è la v.a che conta il numero di operazioni di push verso $v$ al round $t$. Sapendo ora che $\forall v \in V,\ \delta(v) = n$, 
        \[
            \Pr \left( X^{(t)}_{u,v} = 1 \right) = \frac{1}{n} \quad \text{si ha,} \quad
            \mathbb{E}\left[ X^{(t)}_v \right] = \sum_{u \in V_1} \frac{1}{n} = 1
        \]
        
        Applicando ora il Chernoff Bound 
        \[
            \Pr\left( X \geq (1 + \delta)\mu \right) \leq \left( \frac{e^\delta}{(1+\delta)^{(1 + \delta)}} \right)^\mu
        \]
        dove $\mu = 1$, $\delta = c \log n$ con $c > 0$ e tale che $1 + c\log n \geq e^2$ e con $X = X^{(t)_v}$, otteniamo
        \[
            \Pr\left( X^{(t)}_v \geq 1 + c\log n \right) \leq \frac{e^{c\log n}}{(1+ c\log n)^{(1 + c \log n)}} \leq \left(\frac{1}{e}\right)^{c\log n} \leq \frac{1}{n^c}
        \]
    \end{proof}
\end{lemma}

\section{$\mathcal{PULL}$ Broadcast Protocol on Clique}

Consideriamo il solito problema del \textit{broadcast} a singola sorgente $s$ su una clique $K_n$. Definiamo un protocollo di $\mathcal{PULL}$ che risolve tale problema.

Formalemente, possiamo descrivere il problema mediante la tripla $\langle P_{init}, P_{P_final}, G_{pull} \rangle$, dove:
\begin{itemize}
    \item $P_{init}: \exists! x \in V: state(x) = \texttt{ informed } \land \forall y \neq x, state (y) = \texttt{ notInformed}$;
    \item $P_{final}: \forall x \in V: state(x) = \texttt{ informed}$;
    \item $G_{pull}:$ Restrizioni modello Gossip Pull, insieme a $\texttt{KT}$. La rete di comunicazione considerata è una clique e ogni nodo sa di trovarsi in una clique. 
\end{itemize}

Il protocollo $\mathcal{PULL}$ \textit{Broadcast Protocol} (BP) funziona nel seguente modo:

\begin{enumerate}
    \item Inizialmente, la sorgente $s \in V$ è nello stato $\texttt{informed}$ ed è l'unico nodo informato all'interno della rete. Il resto dei nodi non è informato e si trova nello stato $\texttt{notInformed}$.
    \item Ad ogni round $t \leq 1$, ogni nodo $v \in V: state(v) = \texttt{ notInformed}$ esegue un'operazione di \textit{pull} su un'altro nodo $u \in_{u} N(v)$. Se $state(u) = \texttt{informed}$, allora $v$ si fa inviare una copia del messaggio $M$ da $u$ e passa nello stato $\texttt{informed}$.
    \item Il protocollo termina globalmente quando tutti i nodi si trovano dello stato $\texttt{informed}$. 
\end{enumerate}

\begin{teorema}
    Il protocollo BP con input $K_n$ termina in $\theta(\log n)$ w.h.p.
    \begin{proof}
        Fissiamo un time slot $t \leq 0$, e definiamo l'insieme
        \[
            I_0 := \{s\} \qquad \qquad
            I_t := \{v \in V: v \texttt{ is informed at time } t\}
        \]
        e denoteremo con $m_t = |I_t|$.

        La dimostrazione si dividerà in due macro fasi:
        \begin{itemize}
            \item \textbf{Fase 1} in cui si dimostrerà che w.h.p. il numero di nodi \texttt{informed} $|I_t|$ \textit{cresce} in maniera \textbf{esponenziale} fino a $\frac{n}{2}$.
            \item \textbf{Fase 2} in cui si dimostrerà che w.h.p. il numero di nodi \texttt{notInformed} \textit{decresce} in maniera \textbf{esponenziale}.
        \end{itemize}
        
        \begin{description}
            \item[FASE 1] Sia $t_{max} = \max\{t \leq 0: m_t \leq \frac{n}{2}\}$. Fissiamo quindi un tempo $0 \leq t \leq t_{max}$ e definiamo la v.a. binaria $Y_v$ che indica la probabilità che un nodo \texttt{notInformed} diventi \texttt{informed} al tempo $t + 1$.
            \[
                Y_v = Y^{(t+1)}_v = 
                \begin{cases*}
                    1 \texttt{ if } v \texttt{ get informed at time } t + 1 \\
                    0 \texttt{ otherwise}
                \end{cases*}
                \quad \forall v \in V \setminus I_t
            \]
            Sapendo che ogni nodo fa \textit{pull} in maniera \textbf{uniformemente a caso} avremo che
            \[
                \Pr\left( Y_v = 1 | I_t = m_t\right) = \frac{m_t}{n}
            \]
            Possiamo inoltre calcolare la media di tale v.a. che sappiamo sia esattamente uguale alla probabilità che valga 1, quindi
            \[
                \mathbb{E}\left[ Y_v = 1 | I_t = m_t \right] = \frac{m_t}{n}
            \]
            Definiamo ora l'insieme $I^*$ l'insieme dei \textbf{soli nuovi} nodi informati al tempo $t + 1$, ovvero tale che
            \[
                I_{t + 1} = I_t \cup I^*
            \]
            Poiché $I_t$ e $I^*$ sono \textbf{disgiunti} ($I_t \cap I^* = \emptyset$) avremo che $|I_{t+1}| = |I_t| + |I^*|$.
            A noi interessa calcolare il valore atteso del numoro di nodi informati al istante $t+1$, ossia $\mathbb{E}[|I_{t+1}|]$. Si osserva che $|I^*| = \sum_{v \in V \setminus I_t}Y_v$, quindi applicando il valor atteso e la linearità del valor atteso, otteniamo
            \[
                \mathbb{E}\left[ |I^*| \right] = \sum_{v \in V \setminus I} \mathbb{E}[Y_v | I_t = m_t] = (n-m_t)\frac{m_t}{n} \geq \frac{n}{2}\frac{m_t}{n} = \frac{m_t}{2}
            \]
            Perché in questa fase $m_t \leq \frac{n}{2}$, quindi $n - m_t \geq \frac{n}{2}$.
            
            Quindi, 
            \[
                \mathbb{E}\left[ |I_{t+1}|\ |\ |I_t| = m_t \right] = m_t + \mathbb{E}[|I^*|] \geq \frac{3}{2}m_t
            \]

            Abbiamo quindi ottenuto una relazione ricorsiva che ci dice che il numero di nodi informati al tempo $t + 1$ cresce di un fattore costante rispetto al tempo $t$ (finché $m_t \leq \frac{m}{2}$). Perchiò \textit{srotolando} l'equazione avremo che
            \[
                m_t \geq \frac{3}{2}m_t \geq \left(\frac{3}{2}\right)m_{t-2} \geq \dots \geq \left(\frac{3}{2}\right)^t
            \]
            Sia ora il tempo $\tau = \min \{t \geq 1 | m_t > \frac{n}{2}\}$ il primo istante in cui il numero di nodi \texttt{informed} supera la metà dei nodi. Secondo l'equazione di ricorrenza precedente, avremo che
            \[
                \tau \approx \log_{\frac{3}{2}} n/2 \in \theta(\log n)
            \]
            ovvero che in tempo \textit{logaritmico} almeno la metà dei nodi verrà informata.

            Purtroppo abbiamo solo dimostrato che ciò accade in \textbf{media}, mentre il teorema richiede w.h.p.

            Per dimostrare l'alta probabilità suddividiamo la fase 1 in altre 2 sottofasi, in cui verrà dimostrato l'alta probabilità suddidendo per $1 \leq m_t \leq \alpha \log n$ ed $\alpha \log n < m_t \leq \frac{n}{2}$ per qualche valore di $\alpha$.

            \textbf{Sottofase 1.1 (Bootstrap): } $1 \leq m_t \leq \alpha \log n$
            \begin{lemma}
                Per ogni $\alpha > 0$ esiste una costante $\gamma = \gamma(\alpha)$ sufficientemente grande tale che dopo i primi $\tau_1 = \gamma \log n$ time slots, avremo almeno $m_{\tau_1} \geq \alpha \log n$ nodi \texttt{informed}, w.h.p.
                \begin{proof}
                    Invece di considerare singoli istanti di tempo, consideriamo la finestra temporale che va da $[1, \tau_1]$ e definiamo per ogni $v \in V \setminus \{s\}$ la v.a. binaria 
                    \[
                        Y_v^{(\tau_1)} = 
                        \begin{cases*}
                            1 \texttt{ if } v \texttt{ è informed in un round } t \in [1, \tau_1] \\
                            0 \texttt{ otherwise}
                        \end{cases*}
                    \]
                    Poichè le azioni dei nodi sono \textbf{mutualmente indipendenti}, avremo che 
                    \begin{align*}
                        \Pr(Y^{(\tau_1)}_u = 0) 
                        &= \Pr \left( \bigcap_{t = 1}^{\tau_1} u \text{ doesn't pull the message at time } t \right) \\
                        &= \prod_{t = 1}^{\tau_1} \Pr (\text{ doesn't pull the message at time } t) \\
                        &= \prod_{t = 1}^{\tau_1} \left( 1 - \frac{m_t}{n} \right) \\
                        &\leq \left( 1 - \frac{1}{n} \right)^{\tau_1} \quad m_t \geq 1: \text{ almeno la sorgente è informata} \\
                        &\leq  \left( e^{-1/n} \right)^{\tau_1} = e^{-\gamma\frac{\log n}{n}} \quad (1 - x \leq e^{-x})
                    \end{align*}
                    Sia ora $Y^{(\tau_1)}$ la v.a. che conta il numero di nodi informati al round $\tau_1 + 1$, ossia $m_{\tau_1 + 1} = Y^{(\tau_1)}$, data da 
                    \[
                        Y^{(\tau_1)} = \sum_{v \in V} Y_v^{(\tau_1)}
                    \]
                    Applicando il valore atteso e la sua linearità, otteniamo
                    \begin{align*}
                        \mathbb{E}\left[Y^{(\tau_1)}\right] 
                        &= \sum_{v \in V} \mathbb{E}\left[Y^{(\tau_1)}\right] \\
                        &= \sum_{v \in V} 1 - \Pr(Y^{(\tau_1)}_u = 0) \\
                        &\geq \sum_{v \in V} \left(1 - e^{-\gamma\frac{\log n}{n}}\right) \\
                        &\geq n \frac{-\gamma \log n}{2n} = \frac{-\gamma \log n}{2} \qquad (1 - e^{-x} \geq x/2)
                    \end{align*}
                    Osservando che $Y^{(\tau_1)}$ è la somma di v.a. \textbf{indipendenti} tra loro, è possibile applicare il Chernoff Bound nella sua forma moltiplicativa
                    
                    \[
                        \Pr(X \leq (1 - \delta)\mu) \leq e^{\frac{-\mu\delta^2}{2}}
                    \]
                    Ponendo $\delta = \frac{1}{2}$ e $\mu = \gamma \log n$ otteniamo che la probabilità di avere \textbf{meno} di $\frac{-\gamma \log n}{2}$ nodi informati entro i primi $\tau_1 = \gamma \log n$ slots è al più
                    \[
                        \Pr \left( Y^{(\tau_1)} \leq (1 - 1/2)\gamma\log n \right) \leq e^{-\frac{\gamma \log n}{8}} = n^{-\gamma/8}
                    \]
                    
                    Considerando invece un $\gamma = \gamma(\alpha)$ sufficientemente grande, cioè tale che $\frac{\gamma \log n}{2} \geq \alpha \log n$, si ottiene che $m_{\tau_1} \geq \alpha \log n$ w.h.p.
                    \begin{align*}
                        \Pr \left( Y^{(\tau_1)} \leq \alpha \log n \right) &\leq \Pr \left( Y^{(\tau_1)} \leq \frac{\gamma \log n}{2} \right) \leq n^{-\frac{\gamma(\alpha)}{8}} \\
                        & \implies \Pr\left( Y^{(\tau_1)} \geq \alpha \log n \right) \geq 1 - n^{-\frac{\gamma(\alpha)}{8}}
                    \end{align*}
                \end{proof}
            \end{lemma}
            \textbf{Sottofase 1.2: } $\alpha \log n \leq m_t \leq \frac{n}{2}$
            \begin{lemma}
                Esiste un $\beta$ sufficientemente grande, tale che dopo $\tau_2\beta \log n$ time slots avremo che $m_{\tau_2} \geq \frac{n}{2}$ w.h.p.
                \begin{proof}
                    Fissato un $\tau_2 \geq t \geq \tau_1$, sia $m_t = |I_t|$. Notiamo che $m_t$ non è una v.a. in quanto stiamo considerando un $t$ fissato. Sia quindi la v.a. binaria 
                    \[
                        Y^{t}_v = 
                        \begin{cases*}
                            1 \texttt{ if } v \texttt{ risulta è informed al tempo } t + 1 \\
                            0 \texttt{ otherwise}
                        \end{cases*}
                        \forall v \in V \setminus I_t
                    \]
                    Grazie alla prima parte della dimostrazione sappiamo Chernoff
                    \[
                        \mathbb{E}\left[ |I_{t+1}|\ |\ |I_t = m_t| \right] \geq \frac{3}{2}m_t
                    \]
                    e siccome stiamo considerando un $t \geq \tau_1$, dalla dimostrazione della \textbf{sottofase 1.1} abbiamo che
                    \[
                        \mathbb{E}[m_{t+1}] = \mathbb{E}\left[ |I_{t+1}|\ |\ |I_t = m_t| \right] \geq \frac{3}{2}m_t \geq \alpha \log n 
                    \]
                    Importante osservare che $m_{t+1}$ è una \textbf{variabile aleatoria} che dipende da $m_t$, che è fissata. Poiché $m_{t+1}$ è la somma delel v.a. \textbf{indipendenti} $Y_v^t$, possiamo nuovamente applicare il Chernoff Bound e ottenere
                    \[
                        \Pr\left( m_{t+1} \leq (1 - \delta)\frac{3}{2}m_t \right) \leq e^{-\frac{\delta^2}{2}\frac{3}{2}m_t} \leq e^{-\frac{\delta^2}\alpha \log n} = n^{-c}
                    \]
                    Ponendo $\delta \geq \frac{1}{3}$ avremo che
                    \[
                        \Pr\left( m_{t+1} \leq m_t \right) \leq n^{-c}
                    \]
                    dove l'ultimo bound implica l'alta probabilità della crescita esponenziale di $m_t$, ovvero
                    \[
                        \Pr\left( m_{t+1} \geq m_t \right) \geq 1 - n^{-c}
                    \]
                \end{proof}
            \end{lemma}
            \item[FASE 2] Questa fase è del tutto analoga alla prima, con la differenza che basta dimostrare che la decrescita dei nodi \textbf{non informati} è esponenziale nel tempo.
            
            Fissiamo un tempo $t \geq \tau_2 = \beta \log n$. Dato che nella prima fase sappiamo che $m_{\tau_2} \geq \frac{n}{2}$, avremo che un nodo \texttt{notInformed} al tempo $t$ fa \textit{pull} del messaggio e diventa \texttt{informed} al tempo $t+1$ con probabilità \textbf{almeno} $\frac{1}{2}$.

            Fissiamo $z_t = n - m \geq \frac{n}{2}$ il numero di nodi sopravvissuti al tempo $t$. Mediamente avremo che il numero di nodi sopravvissuti al tempo $t+1$ (ovvero quelli ancora \texttt{notInformed}) saranno
            \[
                z_{t+1} \leq \frac{z_t}{2} \leq \frac{z_\tau}{2} \leq \frac{n}{4}
            \]
            Questa catena si può generalizzare come segue
            \[
                z_t \leq \frac{z_{t-1}}{2} \leq \frac{z_{t-2}}{4} \leq \dots \leq \frac{z_{t-i}}{2^i} \leq \dots \leq \frac{z_{\tau_2}}{2^{t - \tau_2}} \leq \frac{n}{2^{t - \tau_2 + 1}}
            \]

            Ci si chiede ora per quali valori di $t$ il numero di nodi sopravvissuti diventa molto piccolo, diciamo $z_t \leq c \log n$. Certamente se $\frac{n}{2^{t - \tau_2 + 1}} \leq c\log n$ allora anche $z_t \leq c \log n$, quindi basta risolvere la prima disuguaglianza per ottenere la prima.
            \begin{align*}
                \frac{n}{2^{t - \tau_2 + 1}} &= c\log n \iff \\
                2^{t - \tau_2} &= \frac{n}{2c\log n} \iff \\
                2^t2^{-\tau_2} &= \frac{n}{2c\log n} \iff \\
                2^t2^{-\beta \log n} &= \frac{n}{2c\log n} \iff \\
                2^tn^{-\beta'} &= \frac{n}{2c\log n} \quad (a^{\log_b c} = c^{\log_b a}) \iff \\
                2^t &= \frac{n^{\beta' + 1}}{2c\log n} \iff \\
                t &= \log \left( \frac{n^{\beta' + 1}}{2c \log n} \right) \\
                &= \log\left(n^{\beta' + 1}\right) - \log (2c\log n) \\
                &= (\beta' + 1)\log n - \theta(\log \log n) \in \theta(\log n)
            \end{align*}
        \item[Final Step] Infine, calcoliamo la probabilità che una volta rimasti a $O(\log n)$ nodi \texttt{notInformed}, in un solo time slot tutti diventano \texttt{informed}, e vediamo che ciò accade w.h.p.
        
        Per comodità diciamo che sono rimasti $c \log n$ nodi \texttt{notInformed}, per qualche costante $c > 0$. La probabilità che un nodo \texttt{notInformed} faccia \textit{pull} su un nodo \texttt{informed} in questa fase del protocollo è dell'ordine di $1 - \frac{\log n}{n}$, perciò si ha che in media, il numero di nodi che non faranno pull su nodi informati sarà
        \[
            \mathbb{E}[z_{t + 1} | z_t = c\log n] = c \frac{\log^2 n}{n}
        \]
        Applicando la disuguaglianza di Markov, avremo che la probabilità di avere \textbf{almeno} un sopravvissuto sarà molto bassa
        \[
            \Pr(z_{t+1} \geq 1) \leq c \frac{\log^2}{n}
        \]
        \[
            \implies \Pr(z_{t+1} = 0) \geq 1- c \frac{\log^2}{n}
        \]
        \end{description}

    \end{proof}
\end{teorema}

\section{Expanders}
\subsection{$\mathcal{PULL}$ Protocol over $\Delta$-Expanders}
\section{Majority Consensus: 3-Maj Protocol}